\documentclass[aspectratio=169,11pt]{beamer}
\usepackage[utf8]{inputenc}
\usepackage[T1]{fontenc}
\usepackage[french]{babel}
\usepackage{helvet}
\renewcommand{\familydefault}{\sfdefault}
\usepackage{hyperref}
\definecolor{bleu}{HTML}{0000FF}
\definecolor{fond}{HTML}{F4F4EF}
\setbeamercolor{background canvas}{bg=fond}
\setbeamercolor{structure}{fg=bleu}
\setbeamercolor{title}{fg=bleu}
\setbeamercolor{frametitle}{fg=bleu,bg=fond}
\setbeamercolor{normal text}{fg=black}
\setbeamertemplate{navigation symbols}{}
\setbeamertemplate{itemize item}{\textcolor{bleu}{\tiny$\blacksquare$}}
\setbeamertemplate{footline}{\hfill\footnotesize\insertframenumber\hspace{1em}\vskip6pt}
\hypersetup{colorlinks=true,urlcolor=bleu,linkcolor=bleu}
\title{Classifications}
\subtitle{Commissions économiques régionales de l'ONU}
\author{Veille institutionnelle}
\date{2026-11-03}

\begin{document}
\begin{frame}\titlepage\end{frame}

\begin{frame}{Classifications}
\begin{itemize}
\item Découpage régional des Nations unies, norme M49
\item Régions géographiques différentes de celles des commissions
\item Pays membres de plusieurs commissions (États-Unis, Canada, pays d'Asie centrale)
\end{itemize}
\end{frame}

\begin{frame}{Classifications : ressources à suivre}
\small
Veille hebdomadaire.
\begin{itemize}
\item \href{https://unstats.un.org/unsd/methodology/m49/}{Norme M49}\\{\footnotesize Documentation, publication : annuelle}
\end{itemize}
\end{frame}

\end{document}
